\documentclass[11pt, letterpaper]{article}

\usepackage[top=.8in, bottom=.8in, left=0.5in, right=0.5in]{geometry}
\usepackage{amsmath, amssymb, amsthm, mathtools}
\usepackage{longtable}
\usepackage{booktabs}
\usepackage{array}
\usepackage{tabularx}
\usepackage{xcolor}
\usepackage{fancyhdr}
\usepackage{titlesec}
\usepackage{setspace}
\usepackage[T1]{fontenc}
\usepackage{enumitem}
\usepackage{tocloft}
\usepackage{tcolorbox}
\tcbuselibrary{breakable, skins}
\usepackage{listings}

\definecolor{RSblue}{RGB}{25, 60, 120}
\definecolor{RSgray}{RGB}{90, 90, 90}
\definecolor{RSgreen}{RGB}{20, 110, 60}
\definecolor{RSred}{RGB}{160, 30, 30}
\definecolor{RSamber}{RGB}{180, 110, 0}
\definecolor{RSpurple}{RGB}{90, 30, 120}

\newtcolorbox{rsbox}{
  enhanced, breakable,
  colback=RSblue!4, colframe=RSblue!60,
  title={\textbf{RS Ontological Statement}},
  fonttitle=\small\bfseries\color{white},
  attach boxed title to top left={yshift=-2mm, xshift=4mm},
  boxed title style={colback=RSblue},
  left=4pt, right=4pt, top=6pt, bottom=4pt
}

\newtcolorbox{verifybox}{
  enhanced, breakable,
  colback=RSgreen!4, colframe=RSgreen!50,
  title={\textbf{External Verification (Continuous Lens)}},
  fonttitle=\small\bfseries\color{white},
  attach boxed title to top left={yshift=-2mm, xshift=4mm},
  boxed title style={colback=RSgreen!70!black},
  left=4pt, right=4pt, top=6pt, bottom=4pt
}

\newtcolorbox{observebox}{
  enhanced, breakable,
  colback=RSamber!5, colframe=RSamber!60,
  title={\textbf{Structural Observation}},
  fonttitle=\small\bfseries\color{white},
  attach boxed title to top left={yshift=-2mm, xshift=4mm},
  boxed title style={colback=RSamber!80!black},
  left=4pt, right=4pt, top=6pt, bottom=4pt
}

\pagestyle{fancy}
\fancyhf{}
\fancyhead[L]{\small\textcolor{RSgray}{RigbySpace Discrete Dynamics}}
\fancyhead[R]{\small\textcolor{RSgray}{Rigby-Tate Correspondence - Page \thepage}}
\fancyfoot[L]{\small\textcolor{RSgray}{D. Veneziano}}
\fancyfoot[R]{\small\textcolor{RSgray}{Discrete Unification Framework}}
\renewcommand{\headrulewidth}{0.4pt}

\titleformat{\section}
  {\large\bfseries\color{RSblue}}{\thesection}{1em}{}[\titlerule]
\titleformat{\subsection}
  {\normalsize\bfseries\color{RSblue}}{\thesubsection}{1em}{}

\onehalfspacing
\setlength{\parskip}{4pt}
\setlength{\parindent}{0pt}

\newcommand{\Z}{\mathbb{Z}}
\newcommand{\SL}{S_{\mathrm{L}}}
\newcommand{\SLpos}{S_{\mathrm{L}}^{+}}
\newcommand{\aop}{\eta}
\newcommand{\rank}{\rho}
\newcommand{\koppa}{\mathbin{\text{\rm\k{o}}}}
\newcommand{\bplus}{\boxplus}
\newcommand{\Gmg}{G_{\mathrm{mg}}}
\newcommand{\Tcharge}{\Gamma}
\newcommand{\Tvac}{T_{\mathrm{vac}}}

\theoremstyle{definition}
\newtheorem{definition}{Definition}[section]
\newtheorem{axiom}{Axiom}[section]
\newtheorem{lemma}{Lemma}[section]
\newtheorem{theorem}{Theorem}[section]
\newtheorem{corollary}{Corollary}[section]
\newtheorem{proposition}{Proposition}[section]

\lstset{
  language=Python,
  basicstyle=\ttfamily\footnotesize,
  commentstyle=\color{gray},
  keywordstyle=\color{blue},
  stringstyle=\color{red},
  showstringspaces=false,
  breaklines=true,
  frame=single,
  numbers=left,
  numberstyle=\tiny,
  xleftmargin=1em,
  xrightmargin=1em,
}

\begin{document}

\begin{center}
  {\LARGE\bfseries\color{RSblue} Rigorous Derivation of the Rigby-Tate Correspondence via Accumulated Causal Depth\par}
  \vspace{0.5cm}
  {\Large\color{RSgray} A Structural Isomorphism between Rational Dynamics and Elliptic Curve Height\par}
  \vspace{0.5cm}
  {\large D. Veneziano\par}
  \vspace{0.5cm}
  {\normalsize \today\par}
\end{center}

\vspace{1cm}

\begin{abstract}
This document establishes a rigorous structural isomorphism between the sequential barycentric evolution of a rational dynamical system and the scalar multiplication operation on elliptic curves. Heuristic magnitude arguments and continuous asymptotic notations are entirely replaced with strict integer derivations using the RS structural rank function and the accumulated causal depth ledger. It is proven that the accumulated structural interaction required to advance the system by $m$ states scales quadratically as exactly $m(m+1)/2$. This discrete quadratic accumulation provides a native, exact integer analogue for the canonical height on elliptic curves, deriving the property strictly from atomic operators without asserting the conclusion as a premise.
\end{abstract}

\section{Introduction}

In continuous arithmetic geometry, the complexity of operations on elliptic curves is measured by the canonical height. The scalar multiplication of a point yields a canonical height that scales quadratically with the scalar multiplier. Current theory successfully utilizes this continuous mapping, but it relies on real-valued limits and abstract geometric spaces. RigbySpace addresses the disconnect between continuous geometric abstractions and the underlying discrete reality by modeling these properties through pure integer structures. 

This document identifies the exact RS-native mechanism that produces this quadratic scaling law. It demonstrates that the quadratic relationship is not an assumed parameter or a continuous approximation, but rather the deterministic result of accumulating linear structural rank in the structural entropy ledger across sequential barycentric additions. The derivation utilizes only permitted RS primitives, maintaining a strict separation between ontological operations and external verification logic.

\section{Ontology and Definitions}

\begin{rsbox}
\begin{definition}[Explicit Rational Pair]
An Explicit Rational Pair (ERP) is an ordered tuple $(n, d) \in \Z \times (\Z \setminus \{0\})$. The numerator $n$ encodes event magnitude; the denominator $d$ encodes geometric capacity. The state space is $\SL$. ERPs are strictly structural and are never reduced.
\end{definition}
\end{rsbox}

\begin{rsbox}
\begin{definition}[Constructive Remainder]
For an integer $x$ and a positive integer $L > 0$, the constructive remainder $x \koppa L$ is the final value $r$ obtained by iteratively subtracting $L$ from $x$ while $r \ge L$, and iteratively adding $L$ to $x$ while $r < 0$. The operation yields $0 \le r < L$ strictly through integer addition and subtraction.
\end{definition}
\end{rsbox}

\begin{rsbox}
\begin{definition}[Structural Rank]
For a positive integer $m$, the structural rank $\rank(m)$ is the total number of prime factors of $m$ counted with multiplicity. By convention $\rank(1) := 0$. Primality is detected strictly by the constructive remainder: an integer $p > 1$ is prime if and only if $p \koppa k \neq 0$ for all integers $k$ with $1 < k < p$.
\end{definition}
\end{rsbox}

\begin{rsbox}
\begin{definition}[Barycentric Addition]
For two ERPs $(n_1, d_1)$ and $(n_2, d_2)$ in $\SL$, the barycentric addition operator $\bplus$ is defined as $(n_1, d_1) \bplus (n_2, d_2) := (n_1 d_2 + n_2 d_1, d_1 d_2)$. The denominator records the multiplicative intersection of the structural histories.
\end{definition}
\end{rsbox}

\begin{rsbox}
\begin{definition}[Accumulated Causal Depth]
The accumulated causal depth $\Phi(m)$ over $m$ sequential states $\{S_1, S_2, \dots, S_m\}$ is the sum of the structural ranks of their respective denominators $d_j$. Formally, $\Phi(m) := \sum_{j=1}^{m} \rank(d_j)$. This represents the total structural entropy deposited into the ledger across the sequence.
\end{definition}
\end{rsbox}

\section{Barycentric Sequence and Rank Growth}

The Mass Gap Barycentric Sequence is constructed from RS primitives without postulating its seeds. The mass gap integer $\Gmg = T_{\mathrm{vac}} \koppa N = 11 \koppa 3 = 2$ is established in the constant lattice as the irreducible per-cycle phase residual. The minimal non-trivial ERP whose denominator is the mass gap integer is $(1, \Gmg) = (1, 2)$; this is the unique unit-magnitude state at mass gap capacity, and it is designated the mass gap seed:
\[
P_{\mathrm{mg}} := (1,\, \Gmg) = (1,\, 2).
\]
The initial vacuum state $S_0 = (1, 1)$ is the unit active seed: the minimal subluminal ERP satisfying $\Tcharge(1,1) = 1^2 - 1^2 = 0$, serving as the structural zero point of the accumulation. The sequence of states $\{S_j\}_{j \ge 0}$ is then generated by repeated barycentric addition of $P_{\mathrm{mg}}$ to the preceding state:
\[
S_j := S_{j-1} \bplus P_{\mathrm{mg}}, \quad j \ge 1.
\]

\begin{lemma}[Strict Linear Rank Growth]
\begin{rsbox}
For the sequence defined by $S_0 = (1, 1)$ and $S_j = S_{j-1} \bplus P_{\mathrm{mg}}$ for $j \ge 1$, the structural rank of the denominator at step $j$ satisfies exactly $\rank(d_j) = j$.
\end{rsbox}
\begin{proof}
\begin{rsbox}
By the definition of barycentric addition, $d_j = d_{j-1} \cdot \Gmg$. Since $d_0 = 1$, iterative substitution yields $d_j = \Gmg^j$.
The mass gap integer $\Gmg = 2$ is prime. The primality condition requires $p \koppa k \neq 0$ for all integers $k$ with $1 < k < p$; for $p = 2$ the open interval $(1, 2)$ contains no integers, so the condition holds vacuously and $\Gmg = 2$ is prime.
The structural rank $\rank(\Gmg^j)$ counts the prime factors of $\Gmg^j$ with multiplicity. Since $\Gmg$ is a single prime, $\rank(\Gmg^j) = j \cdot \rank(\Gmg) = j \cdot 1 = j$. The equality is exact.
\end{rsbox}
\end{proof}
\end{lemma}

\begin{lemma}[Quadratic Causal Depth Accumulation]
\begin{rsbox}
The accumulated causal depth $\Phi(m)$ for the first $m$ states of the Mass Gap Barycentric Sequence is given by the exact relation $2\Phi(m) = m^2 + m$.
\end{rsbox}
\begin{proof}
\begin{rsbox}
By definition, $\Phi(m) = \sum_{j=1}^{m} \rank(d_j)$.
Applying the Strict Linear Rank Growth lemma yields $\Phi(m) = \sum_{j=1}^{m} j$.
The sum of the first $m$ positive integers evaluates exactly to $m(m+1)/2$.
Multiplying both sides by $2$ yields $2\Phi(m) = m^2 + m$. The derivation relies purely on discrete integer addition and structural definitions.
\end{rsbox}
\end{proof}
\end{lemma}

\section{The Rigby-Tate Isomorphism}

\begin{verifybox}
In continuous arithmetic geometry, the canonical height $\hat{h}$ of a point on an elliptic curve satisfies $\hat{h}(mP) = m^2 \hat{h}(P)$ under scalar multiplication by $m$. The multiplier $m$ therefore governs the arithmetic complexity of the resulting point quadratically. This scaling law is an external continuous-lens result; it is presented here for structural comparison only and carries no ontological weight.
\end{verifybox}

\begin{observebox}
The RS result established below --- that the accumulated causal depth $\Phi(m)$ of the Mass Gap Barycentric Sequence satisfies $2\Phi(m) = m^2 + m$ --- exhibits the same dominant $m^2$ scaling as the canonical height. The Rigby-Tate Correspondence identifies this agreement as a structural parallel: the RS integer lattice natively produces quadratic accumulation through linear rank growth, without invoking any continuous geometric postulate.
\end{observebox}

\begin{theorem}[Quadratic Causal Depth]
\begin{rsbox}
For the Mass Gap Barycentric Sequence generated by repeated barycentric addition of the seed $P_{\mathrm{mg}} = (1, \Gmg)$ from the initial state $S_0 = (1, 1)$, the accumulated causal depth satisfies the exact integer identity
\[
2\Phi(m) = m^2 + m.
\]
The dominant scaling of $\Phi(m)$ with $m$ is therefore quadratic.
\end{rsbox}
\begin{proof}
\begin{rsbox}
By the Strict Linear Rank Growth lemma, $\rank(d_j) = j$ for every $j \ge 1$. The accumulated causal depth is therefore
\[
\Phi(m) = \sum_{j=1}^{m} \rank(d_j) = \sum_{j=1}^{m} j.
\]
The finite sum $\sum_{j=1}^{m} j$ equals $m(m+1)/2$. This is verified by the pairing argument: for each $k \in \{1, \ldots, m\}$, pair $k$ with $m+1-k$; each pair sums to $m+1$ and there are $m/2$ such pairs when $m$ is even, or the central term $m+1/2$ is isolated when $m$ is odd, yielding the same closed form in both cases. Multiplying both sides by $2$ gives $2\Phi(m) = m(m+1) = m^2 + m$. Every operation is integer multiplication and addition; no division, no limit, and no continuous quantity enters the derivation.
\end{rsbox}
\end{proof}
\end{theorem}

\section{Conclusion}

The Rigby-Tate Correspondence has been established on two levels. At the ontological level, the Quadratic Causal Depth theorem proves that the Mass Gap Barycentric Sequence, constructed from the RS-derived seed $P_{\mathrm{mg}} = (1, \Gmg)$ and the unit active seed $S_0 = (1,1)$, accumulates causal depth satisfying the exact integer identity $2\Phi(m) = m^2 + m$. The derivation proceeds entirely by integer addition and structural rank definitions; no continuous quantity, no limit, and no floor or ceiling function enters the ontological chain. At the structural observation level, this quadratic accumulation parallels the $m^2$ scaling of canonical height under scalar multiplication in continuous arithmetic geometry, establishing the correspondence as a native RS structural property rather than an external postulate.

\newpage

\section{Appendix: Python Verification Structure}

The following code implements the exact mathematical constructs defined in the document. It utilizes the constructive remainder operator to determine primality and to calculate structural rank, verifying the exact sequence growth and the quadratic causal depth accumulation.

\begin{lstlisting}[language=Python]
def koppa(x, L):
    """
    Constructive Remainder via repeated subtraction/addition.
    No division operator is used.
    """
    if L <= 0:
        raise ValueError("L must be positive")
    r = x
    while r >= L:
        r -= L
    while r < 0:
        r += L
    return r

def is_prime_koppa(p):
    """
    Detects primality strictly using the koppa operator.
    """
    if p <= 1:
        return False
    k = 2
    while k < p:
        if koppa(p, k) == 0:
            return False
        k += 1
    return True

def calculate_rank(m):
    """
    Computes RS Structural Rank: total prime factors with multiplicity.
    """
    if m <= 1:
        return 0
    total = 0
    p = 2
    # Determine prime factors and count multiplicity
    while p <= m:
        if is_prime_koppa(p):
            while koppa(m, p) == 0:
                total += 1
                # Subtract p iteratively to represent integer division 
                # (optimized here as // for simulation speed, 
                # but ontologically it is repeated subtraction)
                m = m // p 
        p += 1
    return total

def barycentric_add(A, B):
    """
    Barycentric Addition of Explicit Rational Pairs.
    (n1, d1) [+] (n2, d2) = (n1*d2 + n2*d1, d1*d2)
    """
    return (A[0] * B[1] + A[1] * B[0], A[1] * B[1])

def verify_structural_isomorphism(max_m):
    """
    Generates the Mass Gap Barycentric Sequence and verifies
    the linear rank growth and quadratic causal depth.
    """
    S = (1, 1)
    P_mg = (1, 2)
    Phi_m = 0
    
    print("m | State Numerator | State Denom | Rank(d) | Phi(m)")
    print("-" * 55)
    
    for m in range(1, max_m + 1):
        S = barycentric_add(S, P_mg)
        
        # Verify Linear Rank Growth Lemma
        rank_d = calculate_rank(S[1])
        assert rank_d == m, f"Linear rank failure at m={m}"
        
        # Accumulate Causal Depth
        Phi_m += rank_d
        
        # Verify Quadratic Causal Depth Accumulation Lemma
        assert 2 * Phi_m == (m * m) + m, f"Quadratic identity failure at m={m}"
        
        print(f"{m:1} | {S[0]:15} | {S[1]:11} | {rank_d:7} | {Phi_m:6}")
        
    print("-" * 55)
    print("All integer bounds and identities verified.")

if __name__ == "__main__":
    verify_structural_isomorphism(10)
\end{lstlisting}

\end{document}